\begin{table}[htbp]
\centering
\caption{Common-window descriptive statistics for the covered economies}
\label{tab:latam_regional_descriptives}
\small
\resizebox{\textwidth}{!}{%
\begin{tabular}{lrrrrrrr}
\toprule
Group & Workers 2017 & Workers 2022 & Share 2017 & Share 2022 & Income 2017 & Income 2022 & Annual growth \\
& \multicolumn{2}{c}{Millions} & \multicolumn{2}{c}{Percent} & \multicolumn{2}{c}{2017 PPP US dollars} & Percent \\
\midrule
Total & 223.1 & 242.8 & 100.0 & 100.0 & 849 & 842 & -0.2 \\
Low & 106.1 & 101.1 & 47.5 & 41.7 & 547 & 549 & 0.1 \\
Middle & 80.2 & 91.5 & 36.0 & 37.7 & 806 & 746 & -1.5 \\
High & 36.8 & 50.1 & 16.5 & 20.7 & 1,814 & 1,611 & -2.3 \\
\bottomrule
\end{tabular}
}
\begin{minipage}{0.98\textwidth}
\footnotesize
\textit{Note:} The table pools the 14 covered economies using their published worker counts. Income is mean monthly labor income per employed person in 2017 PPP US dollars. Annual growth is the compound annual rate between the 2017-Q4 and 2022-Q4 endpoints. Peru refers to Lima and Callao.
\end{minipage}
\end{table}
